%% text-sharing.tex -- shared closing section (ConTeXt)
\subsubsection{Building and sharing}

In \CONTEXT\ the package is the module \code{t-florilegium}: the file
\code{t-florilegium.mkxl} with its Lua half, \code{t-florilegium.lua}, and nine Lua files
beside them that are the same files the \LATEX\ package uses, unchanged. The Lua half is the
glue to \CONTEXT; it takes the place of the \LATEX\ glue in \code{florilegium.lua} and
borrows the rest of that file. The names of the others begin with \code{florilegium-},
because \TeX\ finds a file by its bare name, the document's folder first, and a generic
name would lose to a file of the author's own. After the prefix come \code{ink} (geometry,
color and PDF output), \code{flora} (the primitives), \code{authored} and \code{grow} (the
two compositions), \code{initial} (the decorated initials), and two themes with the figures
one of them draws from (\code{theme-geometric}, the same plant cut from flat shapes;
\code{theme-math}, figures, symbols and equations cascading down the margin;
\code{figures}). The module comes as one zip in the layout of a \TEX\ tree, the files in
\code{tex/context/third/florilegium} and this text with the other examples in
\code{doc/context/third/florilegium}: unzipped into a tree that \CONTEXT\ searches, it lets
any document load the garden with \type{\usemodule[florilegium]} and choose its own seed,
column and hue, in one run of \type{context} at about a second per page. In the repository
the drawing code is kept once, beside the \LATEX\ package, and a continuous integration job
puts the module together and compiles these pages on every push.

\startremark
Three things are worth recording about the port. A literal written into the page bypasses
\CONTEXT's own transformations, so a primitive drawn with \type{\florilegiumdraw} is given
its position and angle as numbers and applies them in Lua. The garden is the page's
background overlay, which \CONTEXT\ builds when the page body, its header and its footer
are already in \type{\pagebox}: walking that box finds the marks the text left, in the
same run. And text that a theme sets in the margin at an angle is turned inside a form
XObject of its own: rotated text in the page's own stream made the backend lose its place,
and the rest of the page's text vanished.
\stopremark
